\subsection{Knowledge Graphs for Enhanced Agent Reasoning}

Two large-scale knowledge graphs were developed for this study: one PFAS-Specific Knowledge Graph derived from PFAS-related scientific papers, and one broad Material Properties Knowledge Graph derived from full abstracts collected using a curated set of PFAS-related material property keywords. The graphs’ sizes in terms of nodes and edges are listed in Table \ref{tab:knowledge graph-sizes}. In both cases, nodes correspond to entities such as materials, properties, or contextual features, while edges capture the relationships between them. The PFAS-Specific Knowledge Graph is designed for depth, focusing on the specific rationale and performance requirements of PFAS in real-world applications. The broader Material Properties Knowledge Graph, by contrast, provides breadth, organizing knowledge around 51 property categories that characterize PFAS. Together, these graphs provide a dual substrate: one enabling targeted exploration within PFAS-specific literature, and the other allowing agents to traverse adjacent domains for more creative and cross-disciplinary hypothesis generation.

\begin{table}[H]
\centering
\caption{Sizes of the two knowledge graphs used in this study.}
\label{tab:knowledge graph-sizes}
\begin{tabular}{lcc}
\hline
\textbf{Knowledge Graph} & \textbf{Nodes} & \textbf{Edges} \\
\hline
PFAS-Specific            & 144,380        & 458,643 \\
Material Properties      & 315,020        & 773,541 \\
\hline
\end{tabular}
\end{table} 

The knowledge graphs are constructed through a systematic pipeline that transforms unstructured scientific literature into structured graph representations. First, a collected corpus of scientific literature is converted into machine readable markdown files and segmented into chunks, each kept within the context length limits of a deployed LLM. Then, each chunk is processed by this LLM to extract semantic triplets, that is two nodes corresponding to entities and an edge representing the entities' relationship, which forms a local subgraph of triplets. Repeating this process across all chunks produces multiple local subgraphs that are subsequently merged into a single large-scale graph. During merging, nodes identified as sufficiently similar based on textual embeddings are consolidated to avoid duplication. This pipeline facilitates generation of knowledge graphs that capture both the breadth of the source corpora and the interconnections across domains. 

\subsection{Selected Models, Hosting, and Supporting Libraries}

The open source LLM \texttt{meta-llama/Llama-3.3-70B-Instruct} \cite{llama3modelcard} was used for both knowledge graph generation and as the agents' language backbone in the multi-agent framework. The model was hosted locally with \texttt{llama.cpp} \cite{gerganov2023llamacpp}, enabling efficient inference and secure deployment through an OpenAI-compatible API schema. To optimize performance, the server was configured with full GPU acceleration, tensor splitting across multiple devices, and an extended context length of 40,000 tokens to accommodate long scientific prompts. The model was deployed in a quantized Q4 format, reducing memory requirements while preserving strong performance, and was loaded directly from internal infrastructure ensuring all data remained on local servers.

This setup provided a lightweight yet flexible inference interface, supporting scalable integration with both the knowledge graph generation pipelines and the multi-agent framework. For knowledge graph generation, the extended context length and efficient local hosting enabled the model to process large volumes of scientific text and extract entities and relations effectively. In the multi-agent framework, the same setup enabled fast and reliable inference with enough contextual capacity to support task decomposition, grounding, and hypothesis generation based on the retrieved information supplied in the context-enriched prompts.

Several open-source Python packages were employed to implement the knowledge graph construction and retrieval pipeline in this study. For text embeddings, the \texttt{sentence-transformers} library \cite{reimers-2019-sentence-bert} was used to load and apply the open source \texttt{nomic-ai/nomic-embed-text-v1.5} model \cite{nussbaum2024nomic}, which produces semantic vector representations of text for similarity search and retrieval-augmented generation. This model was selected for its strong performance on both short- and long-context embedding benchmarks, as well as its open-source availability, which makes it well suited for large-scale knowledge graph embedding and retrieval.

The \texttt{NetworkX} package \cite{SciPyProceedings_11} was used to manipulate knowledge graph data structures, including managing node and edge attributes. For embedding storage and retrieval, the \texttt{Chroma} library \cite{chroma} was adopted for the vector database. Its persistent client enables local storage of embedding collections and supports efficient similarity queries using cosine distance. The open-source \texttt{AutoGen} framework \cite{wu2023autogenenablingnextgenllm} was used to facilitate the multi-agent pipeline, enabling both automated LLM-driven generation and tool usage, such as querying the knowledge graphs.

\subsection{In-House Tools for Knowledge Graph Retrieval}

Several custom tools were used and developed to enable agents to interact with the generated knowledge graphs beyond simple text retrieval. Our \texttt{GraphWeave} method extracts subgraphs linked to document identifiers, collects connected nodes, and summarizes relationships into structured text for agent interpretation. A core component is the \texttt{keyword-to-node} mapping tool, which identifies the closest matching graph nodes to input keywords via embedding comparisons. 

As discussed above, we utilized the \texttt{GraphWeave} method in two tooled approaches to enable the Hybrid GraphWeave and Creative GraphWeave agents. The Hybrid GraphWeave approach combines evidence from both the raw text corpus stored in the \texttt{Chroma} library \cite{chroma} and the knowledge graph by embedding a user query, retrieving top-matching text chunks by cosine-similarity, and linking these chunks to corresponding subgraphs. This dual process ensures that retrieval is grounded in direct textual evidence while also leveraging relational context from the graph. In contrast, the Creative GraphWeave approach focuses on exploratory reasoning within the knowledge graph itself, mapping query keywords to nodes through the \texttt{keyword-to-node} tool, relaxing directional constraints, and applying pathfinding algorithms to assemble subgraphs that expose previously unseen latent structural connections in the knowledge graph. 

Finally, a suite of heuristic graph traversal tools was implemented to extend reasoning across domains. These include Breadth-First Search (BFS) with Semantic Stop, Depth-First Search (DFS), Shortest Simple Path, and Top-$N$ Shortest Simple Path algorithms. These traversal tools allow for varying the enriched context that agents can leverage to propose more holistic and creative hypotheses.

\subsection{Corpus Development}
The corpus for the PFAS-Specific Knowledge Graph and the Material Properties Knowledge Graph was assembled using searches in the Web of Science Core Collection \cite{clarivate2025}. For the PFAS-Specific Knowledge Graph, the query term was ``Per- and polyfluoroalkyl substances''. This search yielded 4,824 results on February 18, 2025, of which 4,716 articles were successfully retrieved, corresponding to a yield of $\approx 97.8\%$. Articles that could not be obtained were primarily unavailable due to missing or inaccurate metadata. Full-text manuscripts were collected through a combination of publisher-provided application programming interfaces and manual scraping with permission. 

The Material Properties Knowledge Graph was constructed from the union of search results generated using each of the 51 property keywords listed verbatim that are non-italicized in Table~\ref{tab:pfas_keywords} Polymer scientists designated the italicized terms “high purity,” “high purity material,” “low surface energy,” and “low surface energy material” as key PFAS material properties; however, they were not incorporated into the knowledge graph due to the prohibitively high volume of manuscripts retrieved.

In contrast to the PFAS-specific collection, full-text acquisition for the Material Properties Knowledge Graph was limited by the sheer volume of papers, the broad distribution of publishers and their varying access policies, and API restrictions that rendered manual scraping impractical. To address this and maximize both coverage and diversity, we leveraged the Web of Science functionality to retrieve abstracts and metadata in batches of 1,000 records. These abstracts constitute the corpus for the Material Properties Knowledge Graph. The final collection, completed on April 21, 2025, contained 160,495 abstracts; however, the Material Properties Knowledge Graph used for multi-agent experimentation was from a random sample of 63,222 abstracts, representing a practical limit given the computational expense of graph merging, traversal, and information retrieval at this scale. We found this size to be sufficient for our Evaluator keywords to identify matches within the knowledge graph. As a preliminary validation of this multiagent pipeline, we observe that even at this final scale, the model demonstrates creativity and yields distinctive outcomes across graph traversal algorithms. The total number of abstracts originally retrieved for each material property keyword is reported in Table~\ref{tab:pfas_keywords}. Future work could improve computational and retrieval efficiency to enable larger-scale expansion and utility of the knowledge graph.

\subsection{Knowledge Graph Construction Pipeline}

We construct the knowledge graph through an established pipeline \cite{buehler_accelerating_2024} that first extracts graph fragments within each paper and then merges them across papers.

\subsubsection{Intra-Paper Extraction}
The retrieved PDFs are first converted to Markdown (\texttt{.md}) using \texttt{marker-pdf} \cite{marker-github}, which preserves core text structure, including section headers, the positions of tables and figures, and in-text reference labels, therefore facilitating downstream LLM-based knowledge extraction. We then segment the text into sections  to improve contextual resolution while reducing token usage. For each section, we perform multi-pass LLM distillation with structured outputs to produce a \emph{graph fragment}, essentially a list of triples, represented as a set of nodes and directed edges between nodes. The extraction uses the following schema in python:

\begin{verbatim}
    class Node(BaseModel):
        id: str
        type: str
            
    class Edge(BaseModel):
        source: str
        target: str
        relation: str
            
    class KnowledgeGraph(BaseModel):
        nodes: List[Node]
        edges: List[Edge]
\end{verbatim}

Then, since terminology within the same paper is highly correlated, we can directly compose all fragments from the same paper into a single knowledge graph, where the identical nodes will connect fragments within a paper.

\subsubsection{Cross-Paper Alignment and Merge}
To handle semantic variations in word use across our cross-domain corpus, we compute node embeddings using a \emph{nomic} sentence embedding model for all nodes so that we can compare them in a high-dimensional vector space; nodes with small angular distance, that is, high cosine similarity (Eq.~\ref{eq:1}) greater than 0.95, are treated as semantically similar and merged. For two nodes with embeddings $\mathbf{v}_i$ and $\mathbf{v}_j$, their cosine similarity is
\begin{equation}
\operatorname{cos\_sim}(\mathbf{v}_i,\mathbf{v}_j)
= \frac{\mathbf{v}_i^\top \mathbf{v}_j}{\lVert \mathbf{v}_i\rVert \,\lVert \mathbf{v}_j\rVert}.
\label{eq:1}
\end{equation}

Following Eq.~\ref{eq:1}, we merge similar node pairs by keeping the higher-degree node and removing the other. This step is essential at our corpus scale, preserving core information while reducing the memory footprint. Using this merging strategy, we iteratively integrate each paper’s component graph into the core knowledge graph, merging duplicates via cosine similarity, and finally produce the large knowledge graphs used in this study. 